\documentclass[10pt,a4paper]{amsart}
\usepackage{amsmath,amssymb,amsthm,graphicx}
% For commutative diagrams we use Xy-pic
\usepackage[all]{xy}

% We use 2cell for 2-commutative diagrams.
\xyoption{2cell}
% Theorem environments.
%
\theoremstyle{plain}
\newtheorem{theorem}[subsection]{Theorem}
\newtheorem{proposition}[subsection]{Proposition}
\newtheorem{lemma}[subsection]{Lemma}

\theoremstyle{definition}
\newtheorem{definition}[subsection]{Definition}
\newtheorem{example}[subsection]{Example}
\newtheorem{exercise}[subsection]{Exercise}
\newtheorem{situation}[subsection]{Situation}

\theoremstyle{remark}
\newtheorem{remark}[subsection]{Remark}
\newtheorem{remarks}[subsection]{Remarks}

\numberwithin{equation}{subsection}

% Macros
%
\def\lim{\mathop{\mathrm{lim}}\nolimits}
\def\colim{\mathop{\mathrm{colim}}\nolimits}
\def\Spec{\mathop{\mathrm{Spec}}}
\def\Hom{\mathop{\mathrm{Hom}}\nolimits}
\def\Ext{\mathop{\mathrm{Ext}}\nolimits}
\def\SheafHom{\mathop{\mathcal{H}\!\mathit{om}}\nolimits}
\def\SheafExt{\mathop{\mathcal{E}\!\mathit{xt}}\nolimits}
\def\Sch{\mathit{Sch}}
\def\Mor{\mathop{\mathrm{Mor}}\nolimits}
\def\Ob{\mathop{\mathrm{Ob}}\nolimits}
\def\Sh{\mathop{\mathit{Sh}}\nolimits}
\def\NL{\mathop{N\!L}\nolimits}
\def\CH{\mathop{\mathrm{CH}}\nolimits}
\def\proetale{{pro\text{-}\acute{e}tale}}
\def\etale{{\acute{e}tale}}
\def\QCoh{\mathit{QCoh}}
\def\Ker{\mathop{\mathrm{Ker}}}
\def\Im{\mathop{\mathrm{Im}}}
\def\Coker{\mathop{\mathrm{Coker}}}
\def\Coim{\mathop{\mathrm{Coim}}}

% Boxtimes
%
\DeclareMathSymbol{\boxtimes}{\mathbin}{AMSa}{"02}

%
% Macros for moduli stacks/spaces
%
\def\QCohstack{\mathcal{QC}\!\mathit{oh}}
\def\Cohstack{\mathcal{C}\!\mathit{oh}}
\def\Spacesstack{\mathcal{S}\!\mathit{paces}}
\def\Quotfunctor{\mathrm{Quot}}
\def\Hilbfunctor{\mathrm{Hilb}}
\def\Curvesstack{\mathcal{C}\!\mathit{urves}}
\def\Polarizedstack{\mathcal{P}\!\mathit{olarized}}
\def\Complexesstack{\mathcal{C}\!\mathit{omplexes}}
% \Pic is the operator that assigns to X its picard group, usage \Pic(X)
% \Picardstack_{X/B} denotes the Picard stack of X over B
% \Picardfunctor_{X/B} denotes the Picard functor of X over B
\def\Pic{\mathop{\mathrm{Pic}}\nolimits}
\def\Picardstack{\mathcal{P}\!\mathit{ic}}
\def\Picardfunctor{\mathrm{Pic}}
\def\Deformationcategory{\mathcal{D}\!\mathit{ef}}
\usepackage{cleveref}
\crefname{originallemma}{original lemma}{original lemmas}
\Crefname{originallemma}{Original lemma}{Original lemmas}
\makeatletter
\newlabel{lemma-mayer-vietoris}{{0A50}{source}{}{sourcepointer-1}{}}
\newlabel{lemma-mayer-vietoris@cref}{{[originallemma][0][]0A50}{[1][0][]source}{}{}{}}
\makeatother
\title{Relative Mayer-Vietoris sequence for the higher direct images along a morphism of schemes}
\author{The Stacks Project Authors}
\date{stacks-project:etale-cohomology.tex:local-snapshot-20261011; GNU FDL 1.2 or later}
\begin{document}
\maketitle
\noindent\textit{Editorial note.} All delivered mathematical text is the author's own, copied from the pinned Stacks Project chapter file etale-cohomology.tex; no statement or proof step was rewritten, reordered or removed.\par\medskip
\noindent\textit{Editorial note.} The proof is the author's complete argument: the short exactness is reduced to the section-level Mayer-Vietoris sequence, which the author explicitly attributes to the proof of the earlier same-chapter lemma, and the rest is the identification of the restricted complexes as injective resolutions.\par\medskip
\noindent\textit{Editorial note.} The single labelled reference occurring inside the counted spans is not undefined in the assembled document: it is carried in assembly.external\_references as a non-load-bearing navigation pointer to the same chapter, so no stub replaces any load-bearing mathematics.\par\medskip
\noindent\textit{Editorial note.} Cross-chapter restatement: the same label occurs in chapter cohomology for a morphism of ringed spaces and an O\_X-module. That statement is neither a special case nor a generalisation of this one, because this statement concerns abelian sheaves on the small etale site and its proof uses fibre products inside Y\_etale, so the two are distinct propositions and are counted separately; this card counts the etale-site version.\par\medskip
\noindent\textit{Editorial note.} The preamble is the chapter's own native preamble excerpt already used for this source: the Xy-pic setup (native preamble lines 14-18) and the native theorem declarations with the mathematical macros (native preamble lines 154-217).\par\medskip
\noindent\textit{Allowed standard prerequisites (editorial).}\begin{itemize}
\item Higher direct images R\textasciicircum{}p f\_* of an abelian sheaf on the small etale site. Same chapter, the definition of R\textasciicircum{}p f\_* F as the cohomology sheaves of f\_* I\textasciicircum{}bullet for an injective resolution F -> I\textasciicircum{}bullet. Supplies the objects in the asserted long exact sequence and the resolution with which the counted proof starts.
\item Restriction of sheaves to open subschemes and injectivity of the restriction of an injective sheaf. Same chapter, the restriction F|\_U of an abelian sheaf to an open subscheme and the fact that I\textasciicircum{}bullet|\_U is an injective resolution of F|\_U. Supplies the identification of the restricted complexes as injective resolutions, which is what makes the cohomology of the restricted complexes the higher direct images of the restrictions.
\item Cohomology of the display of a short exact sequence of complexes. Same chapter, the long exact cohomology sequence attached to a short exact sequence of complexes of abelian sheaves. Turns the displayed short exact sequence of complexes into the long exact sequence of the statement.
\end{itemize}\medskip
\ifdefined\hypertarget\hypertarget{sourcepointer-1}{}\fi
\noindent\textit{Source reference.} Original lemma 0A50: sources/stacks-etale-cohomology/pinned-original-source.tex, source lines 7176--7176. Same-chapter navigation pointer invoked at proof line 7262, where the author cites the earlier lemma for the short exactness of the sequence of section groups. The pointer only resolves that reference for the reader; the cited lemma is not reproduced inside the counted span, and no bibliography citation is involved.\par\medskip
\begin{lemma}[Relative Mayer-Vietoris]
\label{lemma-relative-mayer-vietoris}
Let $f : X \to Y$ be a morphism of schemes. Suppose that $X = U \cup V$
is a union of two open subschemes. Denote
$a = f|_U : U \to Y$, $b = f|_V : V \to Y$, and
$c = f|_{U \cap V} : U \cap V \to Y$.
For every abelian sheaf $\mathcal{F}$ on $X_\etale$
there exists a long exact sequence
$$
0 \to
f_*\mathcal{F} \to
a_*(\mathcal{F}|_U) \oplus b_*(\mathcal{F}|_V) \to
c_*(\mathcal{F}|_{U \cap V}) \to
R^1f_*\mathcal{F} \to \ldots
$$
on $Y_\etale$.
This long exact sequence is functorial in $\mathcal{F}$.
\end{lemma}

\begin{proof}
Let $\mathcal{F} \to \mathcal{I}^\bullet$ be an injective resolution
of $\mathcal{F}$ on $X_\etale$. We claim that we
get a short exact sequence of complexes
$$
0 \to
f_*\mathcal{I}^\bullet \to
a_*\mathcal{I}^\bullet|_U \oplus b_*\mathcal{I}^\bullet|_V \to
c_*\mathcal{I}^\bullet|_{U \cap V} \to
0.
$$
Namely, for any $W$ in $Y_\etale$, and for any $n \geq 0$ the
corresponding sequence of groups of sections over $W$
$$
0 \to
\mathcal{I}^n(W \times_Y X) \to
\mathcal{I}^n(W \times_Y U)
\oplus \mathcal{I}^n(W \times_Y V) \to
\mathcal{I}^n(W \times_Y (U \cap V)) \to
0
$$
was shown to be short exact in the proof of Lemma \ref{lemma-mayer-vietoris}.
The lemma follows by taking cohomology sheaves and using the fact that
$\mathcal{I}^\bullet|_U$ is an injective resolution of $\mathcal{F}|_U$
and similarly for $\mathcal{I}^\bullet|_V$, $\mathcal{I}^\bullet|_{U \cap V}$.
\end{proof}

\end{document}
